\chapter{Groups, Haar measure, Lie algebras}\label{app:groups}
This appendix summarizes the group theory used for symmetries, crossed products and the geometry of de Sitter space.

\section{Haar measure}
A \emph{locally compact group} $G$ carries a left-invariant measure, the \emph{Haar measure}, unique up to scale. It is two-sided for abelian groups and for compact groups; such groups are called \emph{unimodular}. For $G=\mathbb R$ it is $\dd t$, for $U(1)$ it is $\dd\theta/2\pi$ (normalized), and for a compact group it can be normalized to total mass one, which makes averaging $\int_G\dd g\,U(g)$ a projection onto the invariant vectors.

\begin{example}[a non-unimodular group]
The affine group $x\mapsto ax+b$, $a>0$, has left Haar measure $\dd a\,\dd b/a^2$ and right Haar measure $\dd a\,\dd b/a$. The ratio is the \emph{modular function} of the group. It is the group-theoretic cousin of the modular operator $\Delta$ of \cref{ch:tomita}, and it is why crossed products by non-unimodular groups produce weights rather than traces.
\end{example}
For noncompact $G$ the average $\int_G\dd g\,U(g)$ is not a bounded operator; this is the reason group averaging requires a rigged Hilbert space (\cref{ch:groups}).

\section{Characters and duality}
A \emph{character} of an abelian group $G$ is a continuous homomorphism $G\to U(1)$. The characters form the \emph{dual group} $\hat G$ (Pontryagin duality, $\hat{\hat G}=G$):
\[
\hat{\mathbb R}=\mathbb R,\qquad \hat{U(1)}=\mathbb Z,\qquad \hat{\mathbb Z}=U(1).
\]
The Fourier transform $L^2(G)\to L^2(\hat G)$ is unitary (Plancherel). A crossed product by an abelian group $G$ always carries a dual action of $\hat G$, and crossing again gives back the original algebra up to a type I$_\infty$ factor: this is Takesaki duality (\cref{ch:takesaki}).

\section{Lie groups and Lie algebras}
For a Lie group, the tangent space at the identity is the \emph{Lie algebra} $\mathfrak g$ with bracket $[X,Y]$, and one-parameter subgroups are $\ee^{tX}$. In a unitary representation $U(\ee^{tX})=\ee^{-\ii tQ_X}$, and by Stone's theorem (\cref{ch:spectral}) each $Q_X$ is self-adjoint. Compact directions (rotations) give generators with discrete spectrum; noncompact directions (boosts, dilations) give generators with spectrum on all of $\mathbb R$ in nontrivial unitary representations.

\section{Poincar\'e and de Sitter groups}
The Poincar\'e group $\mathrm{ISO}(1,d)$ of $(d{+}1)$-dimensional Minkowski space has translations $P^\mu$ and Lorentz generators $M^{\mu\nu}$. The \emph{boost} $B$ in the $(x^0,x^1)$-plane generates the one-parameter group preserving the Rindler wedge, and its spectrum is $\mathbb R$ (\cref{ch:rs_bw}). The de Sitter group is $SO(1,d+1)$ of dimension $\tfrac12(d+1)(d+2)$. Time translation in a static patch is a boost of the embedding space $\mathbb R^{1,d+1}$, so the static-patch Hamiltonian $H$ has continuous spectrum on $\mathbb R$ just like $B$, a fact behind the type III$_1$ structure of \cref{ch:static_patch}.

\mnote[Haar 1933]{Haar (1933) proved existence of an invariant measure on separable locally compact groups; von Neumann (1934) proved uniqueness for compact ones and Weil (1940) in general. Non-unimodular groups such as the affine group have different left and right Haar measures.}
